\section{All matching-strong C2 sources in the positive contact chart}
\label{sec:matching-source-entry}

We now remove the bounded-angle and bounded-ratio assumptions from the
weak-event entry lemma, provided the four matching packets are strong.
All source hypotheses refer to the same exact positive rational C2 chart;
this is not an assertion that every canonical C2 source has that chart.
Write $C$ for its event weight and $c$ for its contact half-angle.

\begin{theorem}[A complete matching-strong C2 source domain]
Use the physical C2 source with $d=1$, $\lambda=\beta=0$,
\begin{gather*}
 c,a,b,t>0,\qquad 0<r<1,\qquad h=b-a^2t>0,\\
 g=(1+c^2)b-(1-r+a^2c^2t)>0,\qquad j=1+c^2r-t>0,
\end{gather*}
and the original positivity conditions
$1-c^2h>0$, $b-a^2c^2r>0$, with $n,m,A,D>0$.
Suppose $C,\alpha$ are infinitesimal and
\[
 1/n,\quad1/m,\quad1/A,\quad1/D
 \quad\hbox{are bounded}.
\]
Then every independent projective test and real or complex state has
 tensor-SOS full parent. No bound is assumed on $a,c,b,g$ and no
comparability is assumed between $C$ and $\alpha$.
After at most the exact C2 reflection, its actual source has
\[
 c\text{ bounded},\qquad a,b,t\longrightarrow1,
 \qquad r,h,g,j\longrightarrow0.
\]
The arrows denote coefficient residues.
\end{theorem}
\begin{proof}
Let $\mu=\min(n,m,A,D)$. The previously proved matching budget, applied
with $T=C+\alpha$, makes $1/[\mu(C+\alpha)]$ bounded. Since
$C+\alpha$ is infinitesimal, $1/\mu$ is infinitesimal. Thus all four
matching weights are unbounded.

Normalize the actual contact state to have Euclidean norm one. At the exact contact,
\[
 0=4\sum_i e_i p_i+\frac{H}{D_uD_v},\qquad
 \frac{H}{D_uD_v}\ge-2Z^2\ge-2,
\]
where $p_n=X^2$, $p_A=Y^2$ and the other $p_i$ are the normalized
squared event amplitudes. Thus $0\le p_i\le1/(2e_i)$ for each matching
event. Every matching probability has zero residue; in particular
$X,Y$ have zero residue.
Put
\[
 c_u=(1+u_c^2)^{-1/2},\quad s_u=u_c c_u,\qquad
 c_v=(1+v_c^2)^{-1/2},\quad s_v=v_c c_v.
\]
These are bounded positive coefficient quantities with
$c_u^2+s_u^2=c_v^2+s_v^2=1$. Write
$\sigma=\overline{c_us_v}$ and $\tau=\overline{s_uc_v}$.
The two other vanishing matching probabilities give
\[
 \sigma\overline W-\tau\overline Z=0,\qquad
 \tau\overline W-\sigma\overline Z=0.
\]
At least one of $\overline W,\overline Z$ is nonzero, so
$\sigma^2=\tau^2$. Both are nonnegative, hence $\sigma=\tau$.
If they are positive, both contact half-angles are bounded units with
the same residue. If they vanish, either both contact half-angles are
infinitesimal or both are unbounded. Indeed
$c_us_v=s_uc_v=0$ in residue, together with the two unit-circle
identities, permits precisely these two endpoint alternatives.

The exact reflection sends $(u_c,v_c)$ to $(1/v_c,1/u_c)$.
It therefore reduces the two endpoint cases to the infinitesimal one.
It preserves the positive source domain: for $z=c^2$ its parameters are
\[
 c'=\frac1{ac},\quad a'=a,\quad b'=\frac{a^2}{b},\quad
 r'=\frac{a^2zr}{b},\quad t'=\frac{1-zh}{b},
\]
\[
 h'=\frac{a^2zh}{b},\quad g'=\frac{a^2j}{b},\quad j'=\frac gb.
\]
In particular
\[
 1-r'=\frac{b-a^2zr}{b}>0,\qquad
 1-z'h'=\frac{a^2t}{b}>0,\qquad
 b'-a^2z'r'=\frac{a^2(1-r)}b>0.
\]
Direct substitution in the rational source gives
\[
 \Delta'=\frac{\Delta}{b^2z^2},\quad N'=\frac{a^2N}{b^2},\qquad
 (n',m',C',A',D',\alpha')=(n,D,\alpha,A,m,C).
\]
Consequently every hypothesis is retained. These are pure-coefficient
operations. In testing coordinates the reflection is the full homogeneous
parent identity with factor $u^2v^2$, interpreted in projective charts.

It remains to bound $a$ after this orientation. At an ordinary contact
residue $a=v_c/u_c$ is a bounded unit. In the other case both $c$ and
$ac$ are infinitesimal. Suppose for contradiction that $q=a^2$ is
unbounded. Then $z$ and $qz$ are infinitesimal; also
$0<t<1+zr$ is bounded. The numerator in
\[
 m=\frac{(1+qz)(1-r)(1+zr)(1+zr+qzt)}{2g\Delta}
\]
is bounded. Since $1/m$ is bounded, $g\Delta$ is bounded.
The positive expansion of $\Delta$ in the entry lemma gives
$\Delta\ge1$. Thus $g$ is bounded, and
\[
 b=\frac{1-r+g+qzt}{1+z}
\]
is bounded. Positivity of $h=b-qt$ gives $0<t<b/q$, so $t$ is
infinitesimal and $j=1+zr-t$ has residue one. But then
\[
 D=\frac{(1+z)(b+r)(b-qzr)(1+zr+qzt)}{2qj\Delta}
\]
is infinitesimal: the numerator is bounded, since $0<b-qzr<b$,
whereas $1/q$ is infinitesimal, $j$ is a unit, and $\Delta\ge1$.
This contradicts boundedness of $1/D$.

We have proved that $a,c$ are bounded in the oriented source. Its
weak event $C$ is infinitesimal, so the weak-event entry lemma applies:
$a,b,t$ have residue one and $r,h,g,j$ are infinitesimal. If $c$ is
infinitesimal, the completed small-angle radial theorem gives all
independent tests. If $c$ is a bounded unit, the completed ordinary-angle
radial theorem gives all tests. Reflect back when necessary. Restore the
positive pure coefficient normalization if the original $d$ was not one.
\end{proof}

The conclusion strengthens the earlier comparable-weak-packet corollary
on this exact positive contact chart. It covers unbounded ratios
$\alpha/C$ and $C/\alpha$, while preserving independence of actual and
testing parameters. Within the branch where $C,\alpha$ are both infinitesimal, the remaining
C2 sources have at least one weak matching event or have not yet been
entered into this positive contact chart. The theorem does not remove that source-chart obligation.
